\begin{example}[Sample a model of a propositional formula]

    Consider the task of sampling from the models of a propositional formula $\formulawith$, which can be expressed as a \ComputationActivationNetwork{} as
    \begin{align*}
        \frac{1}{\contraction{\formula}} \cdot \formulawith
        = \frac{1}{\contraction{\formula}} \cdot \contractionof{\bencodingofat{\formula}{\headvariable,\shortcatvariables},\tbasisat{\headvariable}}{\shortcatvariables} \, .
    \end{align*}

    Note that in this case, $\frac{1}{2^{\catorder}} \cdot \bencodingofat{\formula}{\headvariable,\shortcatvariables}$ is an ancilla augmentation of $\frac{1}{\contraction{\formula}} \cdot \formulawith$ with $\headvariable$ as the ancilla variable, which q-sample is prepared by the oracle $\qcbencodingof{\formula}$ acting on the ground state $\onehotmapofat{0_{[\catorder+1]}}{\headvariable,\shortcatvariables}$.
    The acceptance rate is then
    \begin{align*}
        \acceptanceprob
        = \frac{1}{2^{\catorder}} \cdot \contraction{\bencodingofat{\formula}{\headvariable=1,\shortcatvariables}}
        = \frac{\contraction{\formula}}{2^{\catorder}}
    \end{align*}
    and the optimal number of rotations is
    \begin{align*}
        \repnumof{\optsymbol}
        = 2 \cdot \mathrm{arcsin}\left(\sqrt{\frac{\contraction{\formula}}{2^{\catorder}}}\right) \, .
    \end{align*}
    The example of a single rotation is sketched in \figref{fig:formulaAmplificationCircuit}.

    \begin{figure}

        \stantikzpic{

            \drawsmalltensor{0}{9}{$\fbasis$}
            \draw (1,9) -- (4,9);
            \drawhblock{5.5}{9}{1.5}{$\qcbencodingof{\exformula}$}
            \draw (5,6) -- (5,8);
            \drawvariabledot{5}{6}
            \draw (6,2) -- (6,8);
            \drawvariabledot{6}{2}


            \draw[dashed] (7.5,12) -- (7.5,-1) node[below] {\corelabelsize $\qstateofat{0}{\headvariable,\shortcatvariables}$};
            \drawtextnode{9}{11}{\corelabelsize $\reflectionof{\goodstatesymbol}$}
            \draw[dashed] (10.5,11) -- (10.5,0.5);
            \drawtextnode{19.5}{11}{\corelabelsize $\reflectionof{0}$}

            \draw[dashed] (28.5,12) -- (28.5,-1) node[below] {\corelabelsize $\qstateofat{1}{\headvariable,\shortcatvariables}$};

            \draw (7,9) -- (8,9);
            \draw (8,8) rectangle (10,10);
            \drawsmalltensor{9}{9}{\corelabelsize $\paulizsymbol$}

            \begin{scope}[shift={(7,0)}]
                \draw (3,9) -- (4,9);
                \drawhblock{5.5}{9}{1.5}{$\qcbencodingof{\exformula}$}
                \draw (5,6) -- (5,8);
                \drawvariabledot{5}{6}
                \draw (6,2) -- (6,8);
                \drawvariabledot{6}{2}

                \draw (7,9) -- (8.5,9);
                \drawsmalltensor{9.5}{9}{$\paulixsymbol$}
                \draw (10.5,9) -- (11.5,9);
                \drawsmalltensor{12.5}{9}{$\paulizsymbol$}
                \draw (13.5,9) -- (14.5,9);
                \drawsmalltensor{15.5}{9}{$\paulixsymbol$}

                \draw (12,6) -- (12,8);
                \drawvariabledot{12}{6}
                \draw (13,2) -- (13,8);
                \drawvariabledot{13}{2}

                \drawsmalltensor{7.5}{6}{$\hgate$}
                \draw (8.5,6) -- (9.5,6);
                \drawsmalltensor{10.5}{6}{$\paulixsymbol$}
                \draw (11.5,6) -- (13.5,6);
                \drawsmalltensor{14.5}{6}{$\paulixsymbol$}
                \draw (15.5,6) -- (16.5,6);
                \drawsmalltensor{17.5}{6}{$\hgate$}
                \draw (18.5,6) -- (23,6);

                \drawsmalltensor{7.5}{2}{$\hgate$}
                \draw (8.5,2) -- (9.5,2);
                \drawsmalltensor{10.5}{2}{$\paulixsymbol$}
                \draw (11.5,2) -- (13.5,2);
                \drawsmalltensor{14.5}{2}{$\paulixsymbol$}
                \draw (15.5,2) -- (16.5,2);
                \drawsmalltensor{17.5}{2}{$\hgate$}
                \draw (18.5,2) -- (23,2);

            \end{scope}

            \begin{scope}[shift={(21,0)}]
                \draw (2.5,9) -- (4,9);
                \drawhblock{5.5}{9}{1.5}{$\qcbencodingof{\exformula}$}
                \draw (7,9) -- (9,9);

                \draw (5,6) -- (5,8);
                \drawvariabledot{5}{6}
                \draw (6,2) -- (6,8);
                \drawvariabledot{6}{2}

            \end{scope}

            \drawsmalltensor{0}{6}{$\fbasis$}
            \draw (1,6) -- (2,6);
            \drawsmalltensor{3}{6}{$\hgate$}
            \draw (4,6) -- (13.5,6);

            \drawsmalltensor{0}{2}{$\fbasis$}
            \draw (1,2) -- (2,2);
            \drawsmalltensor{3}{2}{$\hgate$}
            \draw (4,2) -- (13.5,2);

            \drawtextnode{-3}{9}{\corelabelsize $\headvariable$}
            \draw[dashed] (-2,7.5) -- (30.5,7.5);
            \drawtextnode{-3}{4}{\corelabelsize $\shortcatvariables$}
            \drawtextnode{0}{4.5}{\corelabelsize $\vdots$}

        }
        \caption{Amplitude amplifying circuit for the ancilla augmentation of a function $\exformula$, for $\repnum=1$ repetitions.
        The initial state $\qstateofat{0}{\headvariable,\shortcatvariables}$ is the q-sample of the by $\headvariable$ augmented uniform distributions among the models of $\exformula$, and prepared by action of the oracle $\qcbencodingof{\exformula}$ on the ground state.
        The amplification is then done by the Grover operator, which consists in the reflection $\reflectionof{\goodstatesymbol}$ across the state $\qstateofat{\goodstatesymbol}{\headvariable,\shortcatvariables}$ followed by a reflection reflection $\reflectionof{0}$ across the state $\qstateofat{0}{\headvariable,\shortcatvariables}$.}
        \label{fig:formulaAmplificationCircuit}
    \end{figure}
\end{example}